\section{Relative APS and superconnection setup}
\label{app:aps}

The geometric completion uses APS-type relative data \cite{aps1975}. In practice this means:
\begin{itemize}[leftmargin=1.5em, topsep=3pt, itemsep=2pt]
    \item one works with a declared pair $(D,D_{\mathrm{ref}})$ rather than a single unanchored operator,
    \item boundary conditions on the seam $\Seam$ are part of the definition of the relative index,
    \item the $\eta$-term keeps track of the residual spectral asymmetry,
    \item the superconnection package is used in the same spirit as the spectral-action tradition \cite{chamseddineConnes1997}, but here always with an explicit reference subtraction.
\end{itemize}

This appendix now serves as the proof motor for the geometric reconstruction statements in the main text: its role is to make the common reference structure explicit enough that the reconstruction theorem and the local spectral-scale branch are not left hanging on a slogan.

\section{Relative reflection positivity with fermions}
\label{app:os-fermions}

\subsection{Normalized quotient form}

The relative Euclidean theory used in the main closure section is not interpreted as a difference of measures. Instead it is organized as the normalized quotient
\[
Z_{\mathrm{rel}}[J,\eta,\bar\eta]
=
\frac{Z_{\mathrm{adm}}[J,\eta,\bar\eta]}{Z_{\mathrm{ref}}[0,0,0]}.
\]
Because the denominator is source independent on the admissible sector, the reference theory acts only as a fixed normalizer and does not interfere with reflection positivity.

\subsection{Bosonic and fermionic positivity on \texorpdfstring{$\Padm$}{Padm}}

The commutation relations
\[
[\Theta,P_\Theta]=0,
\qquad
\Theta\Padm=\Padm\Theta
\]
ensure that reflection preserves the admissible subspace. The bosonic part then follows the standard Markov / reflection-stable Osterwalder--Schrader argument on admissible observables. For fermions, the admissible CAR kernel may be written as a reflected quadratic form; equivalently one may express the same condition as positivity of the associated Pfaffian structure on the admissible sector. Hence for every admissible observable $O$ supported in $\tau>0$,
\[
\langle \Theta O,O\rangle_{\mathrm{rel}}
=
\frac{1}{Z_{\mathrm{ref}}[0,0,0]}
\langle \Theta O,O\rangle_{\mathrm{adm}}
\ge 0.
\]
This appendix section is the mechanism behind \Cref{lem:admissible-reflection-positivity,thm:reflection-positivity-admissible,thm:lorentzian-reconstruction}; it is recorded here so that the main-text reconstruction theorem is not supported only by a one-line axiom call.

\section{Admissible OS reconstruction and asymptotic scattering}
\label{app:os-scattering}

This appendix upgrades the admissible Schwinger, reconstruction, and scattering block from a named
standard interface to a manuscript-level proof package. The point is not to replace standard
machinery by ad hoc notation, but to verify explicitly that the admissible branch satisfies the
relevant hypotheses on $\Padm$.

\subsection{Tempered admissible Schwinger families}

\begin{lemma}[Temperedness and moment bounds for admissible Schwinger functions]
\label{lem:tempered-schwinger-bounds}
On the Euclidean continuation of $\mathcal U_{\mathrm{scat}}$, each truncated admissible Schwinger
distribution $S_n^T$ defines a tempered distribution. Moreover, for every compact $K$ and every
source monomial of total degree $n$, its moments obey polynomial bounds determined by the same
uniform heat-kernel estimates and finite local moments of the admissible constructive measure.
\end{lemma}

\begin{proof}
By \Cref{thm:well-posed-primitive-dynamics}, the admissible Euclidean generator is realized by a
self-adjoint elliptic operator on every blocked volume. Together with the constructive measure of
\Cref{thm:constructive-geometric-measure}, this yields uniform heat-kernel bounds and finite moments
for local source insertions. Differentiating $W_{\mathrm{rel}}$ with respect to the local sources
therefore produces distributions whose derivatives are controlled by polynomial combinations of the
same bounds. Passing to the admissible projective limit preserves these polynomial estimates, so the
resulting Schwinger family is tempered.
\end{proof}

\begin{lemma}[Graded Euclidean covariance and symmetry on \texorpdfstring{$\Padm$}{Padm}]
\label{lem:graded-euclidean-covariance-padm}
On the Euclidean continuation of $\mathcal U_{\mathrm{scat}}$, the admissible Schwinger family is
graded Euclidean covariant, graded symmetric under exchange of external legs, and continuous in the
Schwartz topology.
\end{lemma}

\begin{proof}
The Dirac-type principal symbol of the geometric anchor and the low-curvature reduction imply local
Euclidean covariance of the admissible source functional on $\mathcal U_{\mathrm{scat}}$. Bosonic
source derivatives commute, while fermionic derivatives anticommute, so the mixed Schwinger family
is graded symmetric with the usual CAR signs. Continuity follows from the tempered moment bounds of
\Cref{lem:tempered-schwinger-bounds}, which control the action of $S_n^T$ on Schwartz test
functions.
\end{proof}

\subsection{OS and CAR reconstruction on \texorpdfstring{$\Padm$}{Padm}}

\begin{theorem}[Tailored OS/CAR reconstruction on \texorpdfstring{$\Padm$}{Padm}]
\label{thm:tailored-os-car-reconstruction}
Let $\{S_n^T\}_{n\ge 1}$ be the admissible Schwinger family on the Euclidean continuation of
$\mathcal U_{\mathrm{scat}}$. If reflection positivity, temperedness, graded Euclidean covariance,
graded symmetry, and clustering hold on $\Padm$, then there exist:
\begin{enumerate}[label=(\roman*), leftmargin=1.8em, topsep=3pt, itemsep=2pt]
    \item a Hilbert space $\mathcal H_{\mathrm{OS}}$ with vacuum vector $\Omega$,
    \item a strongly continuous positive-energy semigroup $T_t=e^{-tH_{\mathrm{adm}}}$ with
    self-adjoint generator $H_{\mathrm{adm}}\ge 0$,
    \item a reconstructed graded local field net on $\Padm$ whose Euclidean boundary values are the
    original admissible Schwinger functions.
\end{enumerate}
Equivalently, the admissible Euclidean family reconstructs a Wightman/CAR theory on the admissible
sector.
\end{theorem}

\begin{proof}
Define the OS sesquilinear form by
\[
(F,G)_{\mathrm{OS}}:=\langle \Theta F,G\rangle_{\mathrm{rel}}
\]
for positive-time admissible observables. Reflection positivity from
\Cref{thm:reflection-positivity-admissible} and the fermionic quotient-form analysis of
\Cref{app:os-fermions} show that this form is positive semidefinite on the full graded admissible
algebra. Quotienting by the null space and completing gives $\mathcal H_{\mathrm{OS}}$.

Positive Euclidean time translation preserves the null space and acts by contractions, hence defines
a strongly continuous contraction semigroup $T_t$ on $\mathcal H_{\mathrm{OS}}$. The dense domain of
positive-time polynomial observables gives the symmetric generator domain. By Hille--Yosida, the
closure of this generator is self-adjoint and nonnegative, so $T_t=e^{-tH_{\mathrm{adm}}}$ for a
self-adjoint $H_{\mathrm{adm}}\ge 0$.

The temperedness, covariance, graded symmetry, and clustering established in
\Cref{lem:tempered-schwinger-bounds,lem:graded-euclidean-covariance-padm} are precisely the
remaining OS/CAR reconstruction hypotheses. Therefore analytic continuation yields the reconstructed
Wightman/CAR family with vacuum $\Omega$. Because $\Theta$, $P_\Theta$, and $\Padm$ commute, the
entire reconstruction stays on the admissible sector $\Padm$.
\end{proof}

\subsection{Stable massive Haag--Ruelle and LSZ closure}

\begin{lemma}[One-particle spectral projectors from isolated simple poles]
\label{lem:one-particle-projectors-from-poles}
Let $\widetilde G_a^{(2)}(p)$ have an isolated simple pole with positive residue at
$p^2=-m_a^2<0$. Then the corresponding one-particle contribution defines a positive spectral
projector on the admissible Hilbert space.
\end{lemma}

\begin{proof}
The isolated simple pole separates a one-particle mass shell from the multiparticle continuum. By
positivity of the residue, the contour integral of the resolvent around that pole defines a positive
projector on $\mathcal H_{\mathrm{adm}}$. Its range is the one-particle subspace associated with the
stable species $a$.
\end{proof}

\begin{theorem}[Haag--Ruelle scattering on the stable massive admissible sector]
\label{thm:haag-ruelle-admissible-massive}
For every stable massive admissible species with isolated simple pole below threshold, the
Haag--Ruelle asymptotic fields exist on $\mathcal H_{\mathrm{adm}}$ and generate wave operators on
the stable massive scattering sector.
\end{theorem}

\begin{proof}
By \Cref{thm:tailored-os-car-reconstruction,thm:admissible-local-minkowski-net}, the admissible
branch carries a local relativistic net with positive Hamiltonian. By
\Cref{thm:nonperturbative-admissible-qft}, connected correlators cluster exponentially on the
massive sector. The pole assumption and \Cref{lem:one-particle-projectors-from-poles} identify the
one-particle subspace and give positive single-particle spectral projectors. The standard
Haag--Ruelle estimates therefore apply to admissible local operators, yielding the in/out fields and
the wave operators on the stable massive sector.
\end{proof}

\begin{corollary}[LSZ reduction on the stable massive admissible sector]
\label{cor:lsz-admissible-massive}
On the stable massive admissible sector, matrix elements of the scattering operator are obtained by
amputating the connected correlators with the corresponding residue factors.
\end{corollary}

\begin{proof}
By \Cref{thm:haag-ruelle-admissible-massive}, the asymptotic fields and wave operators exist on the
stable massive sector. Applying the usual LSZ reduction to those asymptotic fields and the
admissible local net yields the displayed amputated-correlator formula.
\end{proof}

\subsection{Dressed massless interface in the low-curvature region}

The massless sector is not treated by bare Fock asymptotics in the present manuscript. The relevant
object is the infrared-dressed interface package
\[
\mathcal I_{\mathrm{dress}}
:=
\bigl(\Omega_+^{\mathrm{dress}},\Omega_-^{\mathrm{dress}},S_{\mathrm{dress}}\bigr),
\qquad
S_{\mathrm{dress}}:=(\Omega_+^{\mathrm{dress}})^\ast\Omega_-^{\mathrm{dress}},
\]
constructed from the soft charges of the admissible local net on the low-curvature region
$\mathcal U_{\mathrm{scat}}$. This is exactly the interface package used in the main-text theorem
\Cref{thm:dressed-massless-interface}.

\section{FRW reduction and cosmology interface proofs}
\label{app:frw-interface-proofs}

This appendix section upgrades the cosmology interface block from a proof sketch to a manuscript
proof. The point is to keep the theorem-level content at the level of closed-branch interface data,
while leaving late-time numerical continuation in the comparison layer.

\begin{lemma}[Minisuperspace reduction of the geometric Hodge branch]
\label{lem:minisuperspace-reduction-hodge-branch}
Restrict the closed-branch geometric Hodge action to homogeneous and isotropic variables
\[
(a(t),\varphi(t),\theta_\Sigma(t),N_i(t)).
\]
Then the reduced action is a minisuperspace action whose gravitational and scalar sectors coincide
with the low-curvature FRW/scalaron block used in the main text.
\end{lemma}

\begin{proof}
The closed-branch geometric action is local and covariant. Imposing homogeneity and isotropy kills
all spatial derivative terms and leaves only the FRW scale factor, the scalaron mode, the
determinant-line axion phase, and the heavy-neutrino input package. Evaluating the geometric Hodge
action on this ansatz gives exactly the reduced minisuperspace action displayed in the main-text FRW
theorem.
\end{proof}

\begin{lemma}[Seam-transfer determinant as the infrared vacuum term]
\label{lem:seam-transfer-vacuum-term}
On the closed branch, the positive trace-class seam-transfer operator contributes to the FRW
reduction through the infrared vacuum term
\[
\Lambda_{\mathrm{IR}}
=
\Mplunred^4\bigl[-\log\det\nolimits_{\mathrm{adm}}(1-U_\Sigma)\bigr].
\]
\end{lemma}

\begin{proof}
By \Cref{thm:cosmology-closure}, the seam-transfer operator is positive trace class with spectral
radius strictly below one. Therefore its Fredholm determinant is well defined and contributes an
additive vacuum term to the reduced action. The normalization by $\Mplunred^4$ is exactly the
closed-branch dimensional lift used in the main text.
\end{proof}

\begin{lemma}[Determinant-line axion reduction]
\label{lem:determinant-line-axion-reduction}
On the canonical branch, the determinant-line phase reduces to a compact axion variable with
\[
N_{\mathrm{DW}}=1,
\qquad
\theta_i=\pi\bigl(1-\phiseam(\alpha_\star)\bigr),
\]
and the reheating input is fixed by the scalaron decay width.
\end{lemma}

\begin{proof}
\Cref{thm:boundary-winding-control} fixes the domain-wall number $N_{\mathrm{DW}}=1$ on the
canonical branch. The determinant-line theorem and the electromagnetic closure determine the initial
angle $\theta_i$, while the scalaron branch fixes the decay-width input entering the reheating scale
$T_R$. These are exactly the determinant-line and scalaron data needed in the reduced FRW action.
\end{proof}

\begin{theorem}[Full FRW reduction and cosmology interface theorem]
\label{thm:full-frw-interface-theorem}
The closed-branch minisuperspace reduction determines the cosmology interface package
\[
\bigl(\Lambda_{\mathrm{IR}},N_{\mathrm{DW}},\theta_i,T_R,\mathcal I_{\mathrm{LG}}\bigr)
\]
from the same geometric Hodge, seam-transfer, determinant-line, and neutrino data that determine
the closed branch.
\end{theorem}

\begin{proof}
\Cref{lem:minisuperspace-reduction-hodge-branch} gives the reduced FRW/scalaron action,
\Cref{lem:seam-transfer-vacuum-term} identifies the infrared vacuum term, and
\Cref{lem:determinant-line-axion-reduction} identifies the determinant-line axion and reheating
data. The leptogenesis input block $\mathcal I_{\mathrm{LG}}$ is fixed by the same closed-branch
heavy-neutrino spectrum, family frame, and reheating input used in the main-text leptogenesis
interface theorem. Therefore the entire displayed interface tuple is fixed internally on the closed
branch.
\end{proof}

\section{Yukawa kernels and positivity lemmas}
\label{app:yukawa}

This appendix is the proof motor for the transport positivity, determinant-phase suppression, and
bounded-residual statements used in the main text. The transport program is based on the regularized kernel
\[
\mathcal Y_f^{(\varepsilon)}=(D_f^\dagger D_f+\varepsilon^2)^{-1}.
\]
Three remarks are important.
\begin{enumerate}[leftmargin=1.5em, topsep=3pt, itemsep=3pt]
    \item The regulator $\varepsilon$ is not optional: without it the kernel becomes ill-defined exactly where the admissibility question is most delicate.
    \item Positivity of the kernel belongs to the admissibility-closure layer rather than to the hard carrier kernel. The main text closes it through the finite hexagon gap and the lift to the full admissibility-projected operator; this appendix records the transport-side mechanism behind that closure.
    \item Residual prefactors $\Lambda_{f,j}$ are bounded holonomy data inside the transport theorem. They must remain visible and must not be rhetorically hidden as if the mass map were diagonal without residual transport structure.
\end{enumerate}

\begin{remark}[Positive-kernel determinant route to strong CP]
If the physical mass map is written in the polar form
\[
M_f=U_{f,L}H_fU_{f,R}^\dagger,
\qquad
H_f>0,
\qquad
U_{f,L},U_{f,R}\in SU(3)_F,
\]
with $H_f$ induced by the positive transport kernel $\mathcal Y_f^{(\varepsilon)}$, then the induced determinant phase vanishes:
\[
\arg\det M_f=0.
\]
At the same time weak CP may survive through
\[
V_{\mathrm{CKM}}=U_{u,L}^\dagger U_{d,L}.
\]
Combined with the sheet-CP symmetry protection of $\theta_{\mathrm{sheet}}$ in the relative action,
this is the algebraic core of the determinant-phase theorem stated in the main text and of the
strong-CP closure theorem recorded there. The role of the appendix remark is to display the
polar-decomposition mechanics behind that route, not to reopen its status.
\end{remark}

\subsection{Optional strong-coupling sharpening}

As an appendix-level dynamical sharpening of the hadronic admissibility theorem, one may further require the
admissible center-neutral Wilson functional to obey strong-coupling dominance. On that optional
route the relative Wilson loop satisfies
\[
\langle W(C)\rangle_{\mathrm{rel}}
\le
\exp[-\sigmaQCD A_{\min}(C)],
\qquad
\sigmaQCD=\cthree^2\chi_{\QCD}^2,
\]
and the physical Hilbert space contains no color-nonsinglet asymptotic states. This statement is
kept in the appendix so that the main theorem surface does not depend on an extra strong-coupling
sharpening beyond the closed admissible hadronic sector.

\section{Information-theoretic and cut-and-project readings}
\label{app:info-readings}

This appendix collects interpretive notes that help organize the algebraic structure of the
boundary-polarized closed form. When a theorem-level statement appears below, it refers back to a
main-text proposition; the appendix itself does not upgrade any status label beyond the declared
theorem / derived-consequence / comparison / appendix discipline.

\subsection{\texorpdfstring{$41$}{41}-chain overdetermination of the canonical branch}

The canonical $41$-chain admits a compact overdetermination reading:

\begin{table}[H]
\centering
\small
\begin{tabularx}{\linewidth}{@{}>{\raggedright\arraybackslash}p{0.26\linewidth}>{\raggedleft\arraybackslash}p{0.12\linewidth}Y@{}}
\toprule
Layer & Value & Role \\
\midrule
Electroweak trace block $48\,I_1^{\mathrm{EW}}$ & $41$ & retained electroweak trace packet \\
Block exponent $32\,k_{\mathrm{EW}}$ & $41$ & E$_8$-side compression datum \\
UV hypercharge coefficient $10\,b_1$ & $41$ & one-loop electroweak closure datum \\
Occupancy-weighted carrier norm $\Oadm\gamma+1$ & $41$ & carrier / family compression identity \\
\bottomrule
\end{tabularx}
\caption{$41$-chain overdetermination of the canonical branch. The common integer is read as a
structural consistency marker, not as a cryptographic statement.}
\label{tab:checksum-41}
\end{table}

The value of this table is structural rather than rhetorical: four independently defined layers
collapse onto one integer on the retained branch. This is exactly the sense in which the branch
carries the closed $41$-chain marker $\mathcal I_{41}$.

